\section{Motivación}

Vivimos en una era donde la tecnología está presente en casi todos los ámbitos de nuestra vida, siendo de utilidad tanto para actividades cotidianas como para tareas que presentan dificultades, ya sea porque pueden atentar contra nuestra salud o por limitaciones físicas, entre otras. Gracias a los avances que se han logrado y se siguen logrando, es posible dar soluciones que permitan y faciliten la realización de dichas tareas, sin importar el obstáculo que presenten, mediante sistemas compactos y sencillos de utilizar.

Un ejemplo claro y directo de esta necesidad es la manipulación de materiales peligrosos. En escenarios que involucran sustancias tóxicas, reactivos químicos volátiles, material radiactivo o entornos con temperaturas extremas, la presencia humana es un riesgo inaceptable. Es una de muchas situaciones en las que se requiere la ayuda de la tecnología y, en específico, del sector al que apunta este proyecto: la telerobótica, replicando los movimientos de una mano humana en una mano impresa en 3D. Nuestra motivación principal es demostrar cómo la réplica del movimiento humano en una máquina a distancia puede ofrecer soluciones vitales en situaciones que presentan un alto riesgo o dificultad para las personas.
